\section{Getting Started}
\label{sec:getting-started}

The S\textsuperscript{4}D web interface provides an intuitive platform for analyzing website similarity and structure. This guide walks through all major features and functionality.

\subsection{System Requirements}
\label{subsec:system-requirements}

\begin{itemize}
    \item Modern web browser (Chrome, Firefox, Safari, Edge)
    \item JavaScript enabled
    \item Internet connection for URL processing
    \item Minimum screen resolution: 1024x768
\end{itemize}

\section{Interface Overview}
\label{sec:interface-overview}

The S\textsuperscript{4}D interface consists of four main sections:

\begin{enumerate}
    \item \textbf{Home}: Main URL processing and analysis
    \item \textbf{Experiments}: Batch experiment configuration and execution
    \item \textbf{Results Archive}: Access to previous experiment results
    \item \textbf{Settings}: Theme and configuration options
\end{enumerate}

% TODO: Add screenshot of main interface
% \begin{figure}[htbp]
%     \centering
%     \includegraphics[width=0.9\textwidth]{figures/main_interface.png}
%     \caption{S4D Main Interface Overview}
%     \label{fig:main-interface}
% \end{figure}

\section{Basic URL Analysis}
\label{sec:basic-url-analysis}

\subsection{Adding URLs}
\label{subsec:adding-urls}

\begin{enumerate}
    \item Navigate to the \textbf{Home} section
    \item In the URL input field, enter or paste URLs
    \item \textbf{Multiple URL support}:
        \begin{itemize}
            \item Paste multiple URLs separated by newlines
            \item Paste space-separated URLs
            \item Paste quoted URLs from spreadsheets
        \end{itemize}
    \item Press \textbf{Enter} or click \textbf{Add} to process URLs
    \item Invalid URLs are automatically filtered and reported
\end{enumerate}

\subsection{Selecting Analysis Levels}
\label{subsec:selecting-levels}

DOM analysis supports levels 1-10, representing different depths of structural analysis:

\begin{itemize}
    \item \textbf{Level 1}: Surface-level structural elements
    \item \textbf{Level 5}: Moderate depth analysis (recommended for most cases)
    \item \textbf{Level 10}: Deep structural analysis (resource-intensive)
\end{itemize}

\textbf{Selection Process}:
\begin{enumerate}
    \item Click desired level chips in the Levels section
    \item Selected levels appear highlighted in blue
    \item Multiple levels can be selected for comparative analysis
    \item At least one level must be selected before running analysis
\end{enumerate}

\subsection{Running Analysis}
\label{subsec:running-analysis}

\begin{enumerate}
    \item Ensure URLs are added and levels are selected
    \item Click the \textbf{Run} button
    \item Monitor progress indicator
    \item Use \textbf{Cancel} to abort if needed
    \item Processing time varies based on:
        \begin{itemize}
            \item Number of URLs
            \item Selected levels
            \item Website complexity
            \item Network speed
        \end{itemize}
\end{enumerate}

\section{Understanding Results}
\label{sec:understanding-results}

\subsection{Site Similarity Analysis}
\label{subsec:similarity-analysis}

Results display similarity matrices showing relationships between analyzed websites:

\begin{itemize}
    \item \textbf{Color coding}: Viridis colormap from purple (low similarity) to yellow (high similarity)
    \item \textbf{Percentage values}: Cosine similarity scores (0-100\%)
    \item \textbf{Diagonal elements}: Self-similarity (always 100\%)
    \item \textbf{Interactive features}:
        \begin{itemize}
            \item Hover for detailed tooltips
            \item Click to copy matrix data
            \item PDF export functionality
        \end{itemize}
\end{itemize}

\subsection{Clustering Results}
\label{subsec:clustering-results}

Clustering analysis groups similar website elements:

\begin{itemize}
    \item \textbf{Cluster count}: Number of identified groups
    \item \textbf{DBCV Score}: Quality metric (higher is better)
    \item \textbf{Best level indicator}: Automatically identified optimal analysis level
    \item \textbf{Cluster visualization}: Color-coded HTML outputs
\end{itemize}

\subsection{Generated HTML Files}
\label{subsec:generated-html}

The system generates color-coded HTML visualizations:

\begin{itemize}
    \item \textbf{Organization}: Files grouped by analysis level
    \item \textbf{Naming convention}: level\_domain\_identifier\_guid.html
    \item \textbf{Color coding}: Elements colored by cluster membership
    \item \textbf{Metadata}: URL, GUID, and timestamp information included
\end{itemize}

\section{Advanced Features}
\label{sec:advanced-features}

\subsection{Experiment Configuration}
\label{subsec:experiment-configuration}

For batch processing and research workflows:

\begin{enumerate}
    \item Switch to \textbf{Experiments} tab
    \item Configure experiment parameters:
        \begin{itemize}
            \item Experiment name
            \item URL sets
            \item Level combinations
            \item URLs per domain settings
        \end{itemize}
    \item Launch batch experiments
    \item Monitor progress and results
\end{enumerate}

\subsection{Results Archive}
\label{subsec:results-archive}

Access and analyze previous experiments:

\begin{itemize}
    \item \textbf{Experiment browser}: List of completed analyses
    \item \textbf{Filtering}: Sort by date, success rate, duration
    \item \textbf{Detailed views}: Expand for comprehensive results
    \item \textbf{Data export}: Copy results in multiple formats
    \item \textbf{Generated charts}: View performance visualizations
\end{itemize}

\section{Performance and Optimization}
\label{sec:performance-optimization}

\subsection{Cache System}
\label{subsec:cache-system}

The system implements intelligent caching:

\begin{itemize}
    \item \textbf{URL-level caching}: Previously processed URLs load instantly
    \item \textbf{Cache indicators}: Green indicators show cached content
    \item \textbf{Cache efficiency}: Percentage of reused vs. fresh processing
    \item \textbf{TTL}: 24-hour cache lifetime for CLI, 1-hour for API
\end{itemize}

\subsection{Performance Monitoring}
\label{subsec:performance-monitoring}

Real-time performance metrics:

\begin{itemize}
    \item \textbf{Processing time}: Per URL and per level timing
    \item \textbf{Memory usage}: System resource utilization
    \item \textbf{Success rates}: Processing reliability metrics
    \item \textbf{Bottleneck identification}: Slowest processing steps
\end{itemize}

\section{Troubleshooting}
\label{sec:troubleshooting}

\subsection{Common Issues}
\label{subsec:common-issues}

\textbf{URLs not processing}:
\begin{itemize}
    \item Verify URL format (must include http:// or https://)
    \item Check website accessibility
    \item Ensure stable internet connection
    \item Try reducing number of simultaneous URLs
\end{itemize}

\textbf{Slow processing}:
\begin{itemize}
    \item Reduce analysis levels
    \item Process fewer URLs simultaneously
    \item Check system resources
    \item Utilize cache when available
\end{itemize}

\textbf{Missing results}:
\begin{itemize}
    \item Verify successful processing completion
    \item Check browser console for errors
    \item Refresh page and retry
    \item Contact system administrator if persistent
\end{itemize}

\subsection{Error Messages}
\label{subsec:error-messages}

\begin{itemize}
    \item \textbf{"Invalid URLs skipped"}: Some URLs malformed - check format
    \item \textbf{"No valid URLs found"}: All provided URLs invalid
    \item \textbf{"Processing failed"}: Server error - retry or contact support
    \item \textbf{"Timeout error"}: Request took too long - reduce scope
\end{itemize}

\section{Best Practices}
\label{sec:best-practices}

\subsection{Optimal Usage Patterns}
\label{subsec:optimal-usage}

\begin{itemize}
    \item \textbf{Start small}: Begin with 3-5 URLs to understand results
    \item \textbf{Level selection}: Use levels 1-3 for initial exploration
    \item \textbf{Domain grouping}: Analyze similar domains together
    \item \textbf{Progressive analysis}: Start broad, then focus on interesting clusters
\end{itemize}

\subsection{Result Interpretation}
\label{subsec:result-interpretation}

\begin{itemize}
    \item \textbf{Similarity thresholds}: >80\% indicates high similarity
    \item \textbf{Clustering quality}: DBCV >0.5 indicates good clustering
    \item \textbf{Best level selection}: Use automatically identified optimal levels
    \item \textbf{Visual inspection}: Always review generated HTML files
\end{itemize} 